\documentclass[10pt,letterpaper]{article}
\usepackage{cornell-notes}

\title{Clause 31.11.02.05: Overridden Virtual Functions}
\author{}
\date{}
\setDocTitle{Clause 31.11.02.05: Overridden Virtual Functions}
\setDocSubtitle{C++ 2024}
\setDocOwner{C++ 2024 Cornell Notes}
\setDocDate{\today}
\setCornellCollection{C++ 2024 Cornell Notes}
\setCornellUnitType{Clause}
\setCornellUnitNumber{31.11.02.05}
\setCornellUnitTitle{Overridden Virtual Functions}
\hypersetup{pdftitle={Clause 31.11.02.05: Overridden Virtual Functions},pdfsubject={Cornell notes},pdfauthor={}}

\begin{document}
\maketitle
\begin{CornellOverviewBox}{Overview}
\textbf{Classification:} Standard library - Normative\\[2pt]
\textbf{Learning objective:} Understand the rules, terminology, and semantic consequences associated with Overridden virtual functions.
\end{CornellOverviewBox}\begin{CornellSummaryBox}{Summary}
{\sffamily\bfseries\color{Accent} Summary}\par\vspace{3pt}Section 31.11.2.5 focuses on Overridden virtual functions. Apply the specified interface together with its mandates, effects, result conditions, exception behavior, and complexity constraints. When applying it, preserve the distinction between mandatory requirements, recommendations, permissions, and behavior deliberately left undefined, unspecified, or implementation-defined.
\end{CornellSummaryBox}

\end{document}
